% ==========================================
% Domaine 1 : Chercher
% ==========================================
\newcommand{\CompetenceUn}[1]{
    \ifthenelse{\equal{#1}{1}}
        {Prélever/organiser des informations nécessaires à la résolution de problèmes }{}
    \ifthenelse{\equal{#1}{4}}
        {Sélectionner les informations utiles}{}
}


% ==========================================
% Domaine 2 : Modéliser
% ==========================================
\newcommand{\CompetenceDeux}[1]{
    \ifthenelse{\equal{#1}{1}}
        {Utiliser les mathématiques pour résoudre quelques problèmes \ldots}{}
    \ifthenelse{\equal{#1}{2}}
        {Reconnaître et distinguer des problèmes relevant de situations de proportionnalité \ldots}{}
    \ifthenelse{\equal{#1}{7}}
        {Traduire en langage mathématique une situation réelle en un modèle statistique}{}
}


% ==========================================
% Domaine 3 : Représenter
% ==========================================
\newcommand{\CompetenceTrois}[1]{
    \ifthenelse{\equal{#1}{2}}
        {Produire et utiliser diverses représentations des fractions simples et des nombres décimaux}{}
    \ifthenelse{\equal{#1}{3}}
        {Analyser une figure plane}{}
    \ifthenelse{\equal{#1}{6}}
        {Produire et utiliser des représentations de figures géométriques}{}
    \ifthenelse{\equal{#1}{7}}
        {Reconnaître, analyser et utiliser les éléments de codage d'une figure}{}
    \ifthenelse{\equal{#1}{9}}
        {Produire et utiliser des outils pour représenter un problème (tableaux, graphiques ...)}{}
    \ifthenelse{\equal{#1}{10}}
        {Construire une figure géométrique à l'aide des outils}{}
    \ifthenelse{\equal{#1}{11}}
        {Mesurer et construire des angles avec un rapporteur}{}
}


% ==========================================
% Domaine 4 : Raisonner
% ==========================================
\newcommand{\CompetenceQuatre}[1]{
    \ifthenelse{\equal{#1}{1}}
        {Résoudre des problèmes}{}
    \ifthenelse{\equal{#1}{2}}
        {Raisonner, utiliser des propriétés}{}
    \ifthenelse{\equal{#1}{8}}
        {Résoudre des problèmes faisant appel aux 4 opérations sur les nombres entiers}{}
    \ifthenelse{\equal{#1}{9}}
        {Résoudre des problèmes faisant appel aux 4 opérations sur les nombres décimaux}{}
}


% ==========================================
% Domaine 5 : Calculer
% ==========================================
\newcommand{\CompetenceCinq}[1]{
    \ifthenelse{\equal{#1}{4}}
        {Effectuer des comparaisons, encadrer}{}
    \ifthenelse{\equal{#1}{5}}
        {Calculer avec des nombres entiers}{}
    \ifthenelse{\equal{#1}{6}}
        {Calculer avec des nombres décimaux}{}
    \ifthenelse{\equal{#1}{7}}
        {Calculer un périmètre, une aire ou un volume}{}
    \ifthenelse{\equal{#1}{8}}
        {Calculer, manipuler des horaires et des durées}{}       
}


% ==========================================
% Domaine 6 : Communiquer
% ==========================================
\newcommand{\CompetenceSix}[1]{
    \ifthenelse{\equal{#1}{1}}
        {Utiliser et produire des figures géométriques}{}
    \ifthenelse{\equal{#1}{2}}
        {Fournir le travail attendu à la maison}{}
    \ifthenelse{\equal{#1}{3}}
        {S'impliquer dans le travail à la maison}{}
    \ifthenelse{\equal{#1}{5}}
        {Expliquer une démarche, justifier un choix}{}
    \ifthenelse{\equal{#1}{9}}
        {Utiliser des indicateurs statistiques}{}
}


\newcommand{\CompetenceSept}{
    Connaissances et savoir-faire liés au cours
}

\newcommand{\CompetenceHuit}{
    Travail à la maison réalisé avec sérieux
}

% Les tableaux en haut des sujets
\newcommand{\TableauCompetencesDeux}[2]{
    \begin{center}
        \begin{tabular}{l|c||c||c|c||c|}
            \cline{2-6}
              & I & F & \multicolumn{2}{c||}{S} & TB \\ \hline
            \small{#1} & & & & & \\ \hline
            \small{#2} & & & &  &\\ \hline
        \end{tabular}

        \medskip

        \small{ I : maîtrise \textbf{I}nsuffisante \quad F : maîtrise \textbf{F}ragile \quad S : maîtrise \textbf{S}atisfaisante \quad TB : \textbf{T}rès \textbf{B}onne maîtrise}
    \end{center}
}


\newcommand{\TableauCompetencesTrois}[3]{
    \begin{center}
        \begin{tabular}{l|c||c||c|c||c|}
            \cline{2-6}
             & I & F & \multicolumn{2}{c||}{S} & TB \\ \hline
            \multicolumn{1}{|l||}{\small{#1}} & & & & & \\ \hline
            \multicolumn{1}{|l||}{\small{#2}} & & & & & \\ \hline
            \multicolumn{1}{|l||}{\small{#3}} & & & & & \\ \hline
        \end{tabular}

        \medskip

        \small{ I : maîtrise \textbf{I}nsuffisante \quad F : maîtrise \textbf{F}ragile \quad S : maîtrise \textbf{S}atisfaisante \quad TB : \textbf{T}rès \textbf{B}onne maîtrise}
    \end{center}
}


\newcommand{\TableauCompetencesQuatre}[4]{
    \begin{center}
        \begin{tabular}{l|c||c||c|c||c|}
            \cline{2-6}
             & I & F & \multicolumn{2}{c||}{S} & TB \\ \hline
            \multicolumn{1}{|l||}{\small{#1}} & & & & & \\ \hline
            \multicolumn{1}{|l||}{\small{#2}} & & & & & \\ \hline
            \multicolumn{1}{|l||}{\small{#3}} & & & & & \\ \hline
            \multicolumn{1}{|l||}{\small{#4}} & & & & & \\ \hline
        \end{tabular}

        \medskip

        \small{ I : maîtrise \textbf{I}nsuffisante \quad F : maîtrise \textbf{F}ragile \quad S : maîtrise \textbf{S}atisfaisante \quad TB : \textbf{T}rès \textbf{B}onne maîtrise}
    \end{center}
}